\section{Legal Landscape}
\label{sec:legal}

\subsection{History of the Seats-Votes Curve in Redistricting Litigation}

The seats-votes curve predates all scalar partisan metrics in redistricting
scholarship. Political scientists and legal scholars used the curve to analyze
partisan outcomes of redistricting long before the EG was proposed. The symmetry
standard of \citet{grofmanKing2007} is defined in terms of the $S(v)$ curve:
a map satisfies the symmetry standard if and only if $S(v) = 1 - S(1-v)$ for
all $v$---the curve is anti-symmetric around $(0.50, 0.50)$.

In federal redistricting litigation, the seats-votes curve appeared in expert
testimony in multiple cases preceding \textit{Rucho}, including the North Carolina
and Wisconsin congressional and state legislative redistricting cases in 2016--2018.
Expert witnesses in these proceedings used the curve as a visualization tool for
communicating partisan asymmetry to courts. The district court in \textit{Whitford v.\
Gill}\footnote{\textit{Whitford v.\ Gill}, 218 F.Supp.3d 837 (W.D.\ Wis.\ 2016).}
credited seats-votes curve evidence in finding that Wisconsin's state assembly maps
were a partisan gerrymander.

\textit{Rucho} foreclosed the federal constitutional pathway without addressing the
validity of the seats-votes curve as an analytical tool. The curve's methodological
foundation---the uniform swing model---received no critical attention in \textit{Rucho};
the Court's holding was doctrinal (no judicially manageable standard for federal
courts) rather than methodological.

\subsection{State Court Admissibility Through 2026}

\textbf{North Carolina} (\textit{Harper v.\ Hall}, 2022): The seats-votes curve
was admitted as expert evidence and credited by the NC Supreme Court in its
finding that the enacted congressional map systematically advantaged Republicans.
The court's analysis included reference to the curve's passage through $(0.50,
S(0.50))$ at values substantially below the $(0.50, 0.50)$ fair reference point.

\textbf{Ohio} (\textit{League of Women Voters of Ohio v.\ Ohio Redistricting Comm'n},
2022): Expert testimony including seats-votes curve analysis was credited in the
Ohio Supreme Court's findings that Ohio's congressional map violated the state
constitution's anti-gerrymandering provision. The curve's responsiveness statistics
were part of the evidence demonstrating that Republican-held seats were insulated
from competitive elections.

\textbf{Pennsylvania}: Seats-votes curve analysis was part of the expert testimony
in PA redistricting proceedings. PA courts have focused primarily on comparative
analysis of proposed maps rather than specific metric thresholds, but the curve
has been part of the evidentiary record.

\subsection{Courtroom Advantages}

The seats-votes curve has several advantages over scalar metrics as courtroom evidence:

\textbf{Visual accessibility.} The curve can be plotted on a single graph that
communicates both Partisan Bias and Responsiveness simultaneously without requiring
the court to evaluate multiple separate scalar exhibits.

\textbf{No threshold dependency.} The comparison of bisect and enacted curves
communicates partisan asymmetry visually, without requiring the court to adopt
any specific threshold for EG, MM, or Bias. The court can observe directly that
one curve passes through $(0.50, 0.50)$ and the other does not.

\textbf{Sufficient statistic.} The sufficiency of $S(v)$ for EG and Partisan Bias
(Proposition in Section~\ref{sec:definition}) means that a court admitting the
curve as evidence has admitted the most informative single exhibit possible for
partisan fairness analysis. If the court prefers a scalar summary, EG and Bias
can be read from the curve graph without any additional computation.

\textbf{Temporal robustness.} The uniform swing model can be applied to any election
cycle, and the comparison between bisect and enacted curves is qualitatively
consistent across 2016, 2018, and 2020 election returns. Courts can be shown
that the asymmetry is not an artifact of a single election year.
